\section{Proof complexity}

\begin{result}
\subsection{Cook--Reckhow characterization}
NP equals coNP for every encoded instance type if and only if UNSAT has a
polynomially bounded Cook--Reckhow proof system: a polynomial-time function from
bit strings onto unsatisfiable CNF formulas, with a preimage of polynomial length
for every such formula. Polynomial boundedness asserts that short proofs exist;
it does not give a polynomial-time procedure for finding them.
\citep{CookReckhow1979}
\decl{Complexity/Logics/Propositional/ProofSystem.lean}{320}{Cslib.Complexity.ProofSystem.NP_eq_coNP_iff_exists_polynomiallyBounded}
\end{result}

\begin{result}
\subsection{Resolution size--width tradeoff}
If a CNF formula has initial clause width at most $k$, at most $N$ variables,
and a resolution refutation with at most $s$ lines, it has a refutation of some
integer width $w$ satisfying
\[
w\le k+2\sqrt{2N\ln s}+2.
\]
Thus a large minimum refutation width forces a large refutation size. The
statement concerns proof certificates, not the runtime of a resolution or DPLL
solver.\citep{BenSassonWigderson2001}
\decl{Complexity/Logics/Propositional/Resolution/SizeWidth.lean}{331}{Cslib.Logic.PL.CNF.exists_refutableInWidth_le_of_refutableInSize}
\end{result}
